\begin{abstract}
A identificação de possíveis vínculos entre empresas concorrentes em contratações públicas exige a análise conjunta de relações licitatórias, societárias, de administração e cadastrais, frequentemente dispersas em diferentes bases administrativas. Este trabalho apresenta um \textit{pipeline} computacional determinístico, baseado em grafos heterogêneos multicamada, para apoiar a triagem de sobreposições entre fornecedores que participaram de licitações no município de Divinópolis/MG, no exercício de 2025. Foram utilizados registros das tabelas \texttt{licitacao.csv}, \texttt{habLicitacao.csv}, \texttt{julgLicitacao.csv} e \texttt{quadroSocLicitacao.csv}, provenientes do SICOM/TCE-MG, complementados por informações do Quadro de Sócios e Administradores (QSA) obtidas via BrasilAPI em novembro de 2026. Após o saneamento dos identificadores e a consolidação dos participantes, a rede foi modelada em camadas de participação em licitações, co-disputa por certame e por item, vínculos do QSA e vínculos suspeitos. Para os pares de empresas co-disputantes, calculou-se um escore determinístico de triagem que combina sobreposições do QSA, cadastrais e onomásticas, seguido de filtragem por limiar e redução topológica com remoção de vértices isolados, identificação de componentes conexas e enumeração de motivos triangulares. A análise dos 4.059 pares co-disputantes identificou 202 pares com sobreposição positiva. A redução resultou em um subgrafo com 26 vértices e 23 arestas, distribuído em seis componentes conexas, no qual foram identificados três motivos triangulares em dois componentes. Os demais componentes selecionados apresentaram coincidências cadastrais ou onomásticas sem vínculo direto no QSA. Os resultados constituem indicadores para priorização de verificações documentais, não evidências conclusivas de fraude ou conluio. A integração entre certames e registros de empenhos e pagamentos não foi validada e, por isso, não integra os resultados empíricos deste estudo.
\end{abstract}

\begin{IEEEkeywords}
Teoria dos Grafos, Grafos Heterogêneos Multicamada, Contratações Públicas, Relações Societárias, QSA, SICOM, Triagem de Riscos.
\end{IEEEkeywords}
